\chapter{The four-stage encoder}\label{ch:encoder}

\begin{onething}
The encoder looks at the same 128 $\times$ 128 patch at four zoom levels, from a 32 $\times$ 32 grid of
narrow tokens down to a 4 $\times$ 4 grid of wide ones, and hands all four to the decoder; it hides
tokens at the first zoom level only, and the pyramid itself blurs those holes away.
\end{onething}

You now own every part of one encoder stage: the patch embedding that turns a picture into tokens
(\cref{ch:patchembed}), efficient self-attention (\cref{ch:esa}), and the transformer block that wraps
attention and Mix-FFN with LayerNorms and residuals (\cref{ch:block}). This chapter stacks four of those
stages into a pyramid. The new mathematics is modest: one output-size formula applied four times, a
receptive-field recurrence, and an expected-value calculation that tells you how many stage-2 tokens
see no visible data at all. The new \emph{ideas} are bigger, because the pyramid is what makes masked
training work.

The source is \code{app/foundation_models/components/seg_former_encoder.py}. It is 261 lines\tagC,
over half of them comments. Read it end to end after this chapter.

\section{Why one resolution is not enough}\label{sec:encoder:why}

\subsection{The plain Vision Transformer}

A plain Vision Transformer (ViT) \citep{dosovitskiy2021image} cuts the image into patches once, turns
each patch into a token, and then keeps that one grid of tokens, at one width, through every layer. If
you did that here with 4 $\times$ 4 pixel patches you would have $N = 32 \times 32 = 1024$ tokens in
every layer. Full attention compares every token with every other token, so every layer would build a
score table of
\begin{align*}
N \times N &= 1024 \times 1024 \\
           &= 1{,}048{,}576
\end{align*}
entries per head. Efficient self-attention (\cref{ch:esa}) cuts that down, but the deeper problem is
not cost. It is that one grid fixes one scale of question.

\subsection{Anomalies come in every size}

What the model must learn is ``what normally goes here''. The answer depends on the size of the thing
you are asking about:

\begin{table}[htbp]\centering
\caption{The scale of an anomaly decides which zoom level sees it best.}\label{tab:encoder:sizes}
\begin{tabular}{@{}p{0.36\linewidth}lp{0.3\linewidth}@{}}\toprule
Anomaly & Typical extent & Best seen by \\\midrule
one odd pixel (a flare, a small vessel) & 1 pixel & F1, the only map with that resolution \\
a field of unusual material & about 20 pixels & F2 and F3 \\
a whole region that does not belong & about 200 pixels & F4 \\\bottomrule
\end{tabular}
\end{table}

A single-resolution model must pick one of those rows. A hierarchical encoder returns four feature maps
at four scales and lets the decoder (\cref{ch:decoder}) fuse them, so nothing has to be chosen in
advance.

\begin{analogy}
You are surveying unfamiliar terrain with four maps open at once: a street plan, a city map, a regional
map and a country map. The street plan tells you exactly \emph{where} each building stands but nothing
about the region; the country map tells you \emph{what kind} of region you are in but cannot locate a
house. A surveyor who wants to spot a building that does not belong uses all four together. The
encoder's four stages are those four maps, printed on the same ground: F1 is the street plan, F4 the
country map. The analogy will keep working in \cref{sec:encoder:holes}: if somebody tears holes in the
street plan, the city map still covers the ground underneath.
\end{analogy}

\section{The pyramid, stage by stage}\label{sec:encoder:pyramid}

\subsection{The output-size formula}

Every stage starts with a patch embedding, which is a convolution (\cref{ch:shapes}, \cref{ch:patchembed}).
A reminder of what you need: a \index{convolution!output size}convolution with kernel size $k$, stride
$s$ and padding $p$ turns an input grid of side $H_{\text{in}}$ into an output grid of side
\begin{equation}\label{eq:encoder:convout}
H_{\text{out}} = \left\lfloor \frac{H_{\text{in}} + 2p - k}{s} \right\rfloor + 1 .
\end{equation}
In words: pad both sides by $p$, slide a window of width $k$ in steps of $s$, and count the positions
at which the window still fits. The floor $\lfloor\cdot\rfloor$ rounds down, because a window that would
hang over the edge is not computed.

The encoder uses two settings, read from the constructor (\code{seg_former_encoder.py}, lines 108--123):
\begin{itemize}
\item stage 1: $k = 4$, $s = 4$, $p = 0$ (non-overlapping 4 $\times$ 4 blocks);
\item stages 2--4: $k = 3$, $s = 2$, $p = \lfloor 3/2 \rfloor = 1$ (overlapping windows).
\end{itemize}

\begin{worked}[label={we:encoder:grids}]{The four grid sizes from the formula}
Do this on paper. The input to stage 1 is the compressor's output, 128 $\times$ 128.

Stage 1, $k=4$, $s=4$, $p=0$:
\begin{align*}
H_1 &= \left\lfloor \frac{128 + 2\cdot 0 - 4}{4} \right\rfloor + 1 && \text{substitute}\\
    &= \left\lfloor \frac{124}{4} \right\rfloor + 1 && \text{$128 + 0 - 4 = 124$}\\
    &= \lfloor 31 \rfloor + 1 && \text{$124/4 = 31$ exactly}\\
    &= 32 .
\end{align*}
Stage 2, $k=3$, $s=2$, $p=1$, input 32:
\begin{align*}
H_2 &= \left\lfloor \frac{32 + 2\cdot 1 - 3}{2} \right\rfloor + 1
     = \left\lfloor \frac{31}{2} \right\rfloor + 1
     = \lfloor 15.5 \rfloor + 1
     = 15 + 1 = 16 .
\end{align*}
Stage 3, input 16:
\begin{align*}
H_3 &= \left\lfloor \frac{16 + 2 - 3}{2} \right\rfloor + 1
     = \left\lfloor \frac{15}{2} \right\rfloor + 1
     = 7 + 1 = 8 .
\end{align*}
Stage 4, input 8:
\begin{align*}
H_4 &= \left\lfloor \frac{8 + 2 - 3}{2} \right\rfloor + 1
     = \left\lfloor \frac{7}{2} \right\rfloor + 1
     = 3 + 1 = 4 .
\end{align*}
So the grids are $128 \to 32 \to 16 \to 8 \to 4$. Notice that in stages 2--4 the floor throws away a
half every time: the padded input has an odd length ($31$, $15$, $7$), and the last half-step does not
fit a whole window. That is why $p = 1$ with $k = 3$, $s = 2$ halves an even side exactly.
\end{worked}

Every grid is square, so the token count of stage $i$ is $N_i = H_i \times W_i = H_i^2$:
\begin{align*}
N_1 &= 32 \times 32 = 1024, & N_2 &= 16 \times 16 = 256,\\
N_3 &= 8 \times 8 = 64,     & N_4 &= 4 \times 4 = 16.
\end{align*}
Each stage has a quarter as many tokens as the one before: halving each side quarters the area.

\subsection{Widths go up as grids go down}

Each patch embedding also sets the \index{embedding width}\emph{embedding width} $C_i$, the number of
features every token carries. The configuration is the SegFormer-B0 one \citep{xie2021segformer}:
$C = (32, 64, 160, 256)$. Two quantities move in opposite directions:
\begin{itemize}
\item \textbf{spatial detail falls}, from a 32 $\times$ 32 grid to a 4 $\times$ 4 grid;
\item \textbf{channel width rises}, from 32 to 256 numbers per token.
\end{itemize}
Early stages know \emph{where} things are but little about \emph{what} they are; late stages know a
great deal about \emph{what} but almost nothing about \emph{where}. A pyramid like this, with grid size
falling and width rising stage by stage, is called a \index{feature pyramid}\emph{feature pyramid}.

How much information does each stage hold? Count the numbers in one feature map, $N_i \times C_i$:
\begin{align*}
\text{F1}: \; 1024 \times 32 &= 32{,}768,\\
\text{F2}: \; 256 \times 64 &= 16{,}384,\\
\text{F3}: \; 64 \times 160 &= 10{,}240,\\
\text{F4}: \; 16 \times 256 &= 4{,}096.
\end{align*}
The width grows by $\times 2$, $\times 2.5$, $\times 1.6$ while the token count falls by $\times 4$
each time, so each map holds fewer numbers than the last. The pyramid \emph{summarises}: it is a
compression, not a reshuffle. Compare the input to stage 1: $32 \times 128 \times 128 = 524{,}288$
numbers. F1 already holds sixteen times fewer.

\begin{figure}[htbp]\centering
\begin{tikzpicture}[x=1cm,y=1cm,font=\small]
  % one ground square = 2.6 cm for 128 pixels; channel bars at 2.6 cm per 256 channels
  \def\side{2.6}
  \foreach \i/\n/\c/\lab/\shp [count=\k from 0] in
     {1/32/32/F1/{(B, 32, 32, 32)},
      2/16/64/F2/{(B, 64, 16, 16)},
      3/8/160/F3/{(B, 160, 8, 8)},
      4/4/256/F4/{(B, 256, 4, 4)}}{
    \pgfmathsetmacro{\xo}{\k*3.45}
    \pgfmathsetmacro{\st}{\side/\n}
    \begin{scope}[shift={(\xo,0)}]
      \fill[spat!8] (0,0) rectangle (\side,\side);
      \ifnum\k=0
        % hidden tokens at stage 1 only: a fixed pseudo-random 65% pattern
        \foreach \r in {0,...,31}{\foreach \q in {0,...,31}{
          \pgfmathtruncatemacro{\h}{mod(7*\r+13*\q+3*\r*\q,20)}
          \ifnum\h<13 \fill[maskgrey] (\q*\st,\side-\r*\st-\st) rectangle ++(\st,\st);\fi}}
      \fi
      \draw[spat!70,line width=0.15pt,step=\st] (0,0) grid (\side,\side);
      \draw[spat,thick] (0,0) rectangle (\side,\side);
      \node[above,font=\bfseries\small,text=spat] at (\side/2,\side) {\lab};
      \node[below] at (\side/2,-0.05) {\n\,$\times$\,\n\ = \pgfmathparse{int(\n*\n)}\pgfmathresult\ tokens};
      \pgfmathsetmacro{\bw}{\side*\c/256}
      \fill[spat] (0,-0.95) rectangle (\bw,-0.7);
      \node[right,font=\footnotesize] at (\bw,-0.825) {};
      \node[below,font=\footnotesize] at (\side/2,-1.0) {width $C=\c$};
      \node[below,font=\footnotesize\ttfamily] at (\side/2,-1.45) {\shp};
    \end{scope}}
  \foreach \k/\t in {0/{k4 s4 from input},1/{k3 s2},2/{k3 s2}}{
    \pgfmathsetmacro{\xa}{\k*3.45+\side+0.05}
    \draw[-{Stealth},thick,spat] (\xa,\side/2) -- ++(0.75,0);}
\end{tikzpicture}
\caption{The encoder pyramid drawn to scale. Each square covers the same 128 $\times$ 128 pixels of
ground; its grid lines are the stage's tokens (stage 1: one token per 4 $\times$ 4 pixels, stage 4: one
per 32 $\times$ 32). The bar under each square is the token width, drawn at one scale (256 channels =
the full square width). Grey cells in F1 are hidden tokens at a 65\% hiding rate (an illustrative
pattern); only F1 has them. Arrows are patch embeddings (stride 2, kernel 3).}
\label{fig:encoder:pyramid}
\end{figure}

\coursefig[0.98\linewidth]{iw-16-encoder}{The course's view of the same pyramid: each patch embedding
(arrow) shrinks the grid, by 4 at stage 1 and by 2 after, and widens the channels, producing F1 to F4 at
1/4 to 1/32 of the input resolution. Hidden tokens (hatched) exist at stage 1 only. Front-face size and
depth are compressed for drawing; \cref{fig:encoder:pyramid} is to scale.}{fig:encoder:course}

\subsection{The full configuration}

\begin{table}[htbp]\centering
\caption{The encoder configuration. The last two columns are derived in
\cref{sec:encoder:attncost}.}\label{tab:encoder:config}
\begin{tabular}{@{}crrrrrrrr@{}}\toprule
Stage & $C_i$ & heads & head dim & $R_i$ & blocks & grid & $N_i$ & keys $N_i/R_i^2$ \\\midrule
1 & 32  & 2 & 16 & 8 & 2 & 32 $\times$ 32 & 1024 & 16 \\
2 & 64  & 2 & 32 & 4 & 2 & 16 $\times$ 16 & 256  & 16 \\
3 & 160 & 5 & 32 & 2 & 2 & 8 $\times$ 8   & 64   & 16 \\
4 & 256 & 8 & 32 & 1 & 2 & 4 $\times$ 4   & 16   & 16 \\\bottomrule
\end{tabular}
\end{table}

The \index{head dimension}head dimension is $C_i$ divided by the number of heads (\cref{ch:attention}):
$32/2 = 16$, $64/2 = 32$, $160/5 = 32$, $256/8 = 32$.

\begin{sourcenote}[: the number of heads]
The model class's default is \lstinline|num_heads = [1, 2, 5, 8]|
(\file{app/foundation_models/components/hyperspectral_seg_former_mae.py}, line 71), and the encoder's
docstring table shows 1 head at stage 1 (\file{seg_former_encoder.py}, line 80). The configurations
that trained the shipped models override it to \lstinline|[2, 2, 5, 8]|
(\file{configs/hyperspectral_segformer_exp_2.json}, line 33, and
\file{research/experiments/indradhanu-finetune/FT04-final-joint/config.json}). The v0.1.0
configuration used \lstinline|[1, 2, 5, 8]|. Heads change no parameter count (\cref{ch:params}), so a
checkpoint from either loads into either shape, but they change what the model computes. Trust the
config, not the default.
\end{sourcenote}

\begin{drift}
Configuration values as of 2026-10-29: \lstinline|num_heads = [2, 2, 5, 8]|\tagM{configs/hyperspectral_segformer_exp_2.json},
\lstinline|embed_dims = [32, 64, 160, 256]|, \lstinline|reduction_ratios = [8, 4, 2, 1]|,
\lstinline|num_blocks = [2, 2, 2, 2]|\tagM{research/experiments/indradhanu-finetune/FT04-final-joint/config.json}.
Indradhanu Heavy (\cref{ch:heavy}) widens the stages; everything in this chapter is for the
standard widths.
\end{drift}

\section{The attention budget per stage}\label{sec:encoder:attncost}

Efficient self-attention shrinks the keys and values by a \index{reduction ratio}reduction ratio $R$
in each direction (\cref{ch:esa}). A grid of $H_i \times W_i$ tokens becomes
$(H_i/R_i) \times (W_i/R_i)$ keys, so the key count is $N_i / R_i^2$:
\begin{align*}
\text{stage 1:}\quad \frac{1024}{8^2} &= \frac{1024}{64} = 16,\\
\text{stage 2:}\quad \frac{256}{4^2} &= \frac{256}{16} = 16,\\
\text{stage 3:}\quad \frac{64}{2^2} &= \frac{64}{4} = 16,\\
\text{stage 4:}\quad \frac{16}{1^2} &= 16.
\end{align*}
Every stage attends over exactly 16 keys. That is not a coincidence: the reduction ratios $(8, 4, 2, 1)$
halve as the grid halves, so each stage's keys form the same 4 $\times$ 4 summary of the patch. The
score table per head is queries times keys, $N_i \times 16$:
\[
1024 \times 16 = 16{,}384, \quad 256 \times 16 = 4{,}096, \quad 64 \times 16 = 1{,}024, \quad 16 \times 16 = 256.
\]
Stage 1 has the most queries and therefore the largest table, but the table is $64$ times smaller than
the $1{,}048{,}576$ entries of full attention from \cref{sec:encoder:why}.

\begin{sourcenote}[: masked stage 1 falls back to full attention]
When tokens have been removed, stage 1 holds fewer than $H \times W$ tokens, so ESA cannot fold them
back into a grid for its reduction convolution. It tests \lstinline|can_reduce = (self.reduction_ratio > 1) and (N == H * W)| (\file{app/foundation_models/components/efficient_self_attention.py}, line 144) and,
when that fails, uses the kept tokens themselves as keys: full attention among the survivors. Mix-FFN
likewise skips its depthwise convolution when $N \ne H \times W$
(\file{mix_ffn.py}, lines 118--137). Two comments in \file{seg_former_encoder.py} describe something
else: lines 203--206 say ``ESA's reshape will produce a grid with gaps, which is fine: the reduction
conv pools over them'', and lines 219--223 say ``we pass the actual token grid dims instead''. Neither
happens: the encoder passes the \emph{full} $H, W$, and ESA takes its full-attention fallback. The
source code wins over its comments.

The consequence for cost: with 359 kept tokens (\cref{sec:encoder:walk}) the masked score table is
$359 \times 359 = 128{,}881$ entries per head\tagC, about eight times the unmasked reduced table of
$16{,}384$. Masking makes the linear layers of stage 1 about three times cheaper, but makes its
attention scores more expensive. The ESA comment (line 142) estimates 262K entries ``very cheap'' for
50\% masking; at 65\% it is 129K.
\end{sourcenote}

\section{How far each token sees: the receptive field}\label{sec:encoder:rf}

\begin{definition}[Receptive field]
The \index{receptive field}\emph{receptive field} of a token is the set of input pixels that can
change its value through the convolutions that built it. For a square window we quote its side, in
pixels.
\end{definition}

Attention lets every token read every other token, so in a transformer the receptive field is in
principle the whole patch. But the \emph{patch embeddings} alone already fix how many pixels go into a
token \emph{before} attention runs, and that is what decides how fine a stage's picture is. Two numbers
track it:
\begin{itemize}
\item $r_i$, the receptive field side, in input pixels, of one stage-$i$ token;
\item $j_i$, the \index{jump}\emph{jump}: how many input pixels apart two neighbouring stage-$i$ tokens are.
\end{itemize}

A convolution with kernel $k$ and stride $s$ applied to a grid with jump $j$ and receptive field $r$
gives
\begin{equation}\label{eq:encoder:rf}
r_{\text{out}} = r + (k - 1)\, j, \qquad j_{\text{out}} = j \, s .
\end{equation}
Derive it in words. The output window covers $k$ neighbouring input tokens. The first of them sees $r$
pixels. Each further token is $j$ pixels along, so it adds $j$ new pixels at the far edge. There are
$k - 1$ further tokens, giving $r + (k-1)j$. Neighbouring outputs are $s$ input tokens apart, each $j$
pixels apart, so the new jump is $s \cdot j$.

\begin{worked}[label={we:encoder:rf}]{Receptive fields of the four stages}
Start from the input pixels: $r_0 = 1$ (a pixel sees itself), $j_0 = 1$. Do each line on paper.

Stage 1 ($k = 4$, $s = 4$):
\begin{align*}
r_1 &= r_0 + (4 - 1)\, j_0 = 1 + 3 \cdot 1 = 4, & j_1 &= j_0 \cdot 4 = 4 .
\end{align*}
Stage 2 ($k = 3$, $s = 2$):
\begin{align*}
r_2 &= r_1 + (3 - 1)\, j_1 = 4 + 2 \cdot 4 = 12, & j_2 &= 4 \cdot 2 = 8 .
\end{align*}
Stage 3:
\begin{align*}
r_3 &= 12 + 2 \cdot 8 = 28, & j_3 &= 8 \cdot 2 = 16 .
\end{align*}
Stage 4:
\begin{align*}
r_4 &= 28 + 2 \cdot 16 = 60, & j_4 &= 16 \cdot 2 = 32 .
\end{align*}
A stage-1 token sees exactly its own 4 $\times$ 4 block: no overlap with its neighbours, because $r_1
= j_1$. A stage-4 token is built from a 60 $\times$ 60 window, and its neighbours are 32 pixels apart:
neighbouring windows overlap by $60 - 32 = 28$ pixels. The jumps $4, 8, 16, 32$ are the ``1/4 to 1/32
scale'' of the stage names.
\end{worked}

That $r_1 = j_1$ at stage 1 is a deliberate design choice, and the masking chapter (\cref{ch:masking})
depends on it: a stage-1 token is built only from its own 16 pixels, so hiding a token hides exactly
those pixels and nothing leaks across token borders.

\begin{sourcenote}[: kernel 4, not 7]
The patch-embedding module's docstring describes stage 1 as ``kernel=7, stride=4 -> overlap = 7 - 4 = 3
pixels'' (\file{app/foundation_models/components/overlap_patch_embedding.py}, line 10), as in the
original SegFormer. The encoder builds stage 1 with \lstinline|patch_size = 4 if i == 0 else 3|
(\file{seg_former_encoder.py}, line 110) and \lstinline|padding = 0| (line 114). The source wins: kernel 4,
no overlap. With kernel 7 the receptive field would be $1 + 6 = 7$ pixels and a hidden token's pixels
would leak into its visible neighbours.
\end{sourcenote}

\section{One stage, in code}\label{sec:encoder:loop}

The encoder's \code{forward} is one loop that runs four times:
\begin{lstlisting}
for i in range(self.num_stages):
    x, H, W = self.patch_embeds[i](x)        # -> (B, N, C_i)
    B, N, C = x.shape                        # remember the FULL token count

    gather_indices = None
    if i == 0 and keep_mask is not None:     # hide tokens, stage 1 only
        x, gather_indices = TokenMasking.remove_tokens(x, keep_mask)
                                             # (B, N_kept, C_1)

    for block in self.blocks[i]:             # 2 x SegFormerBlock, shape unchanged
        x = block(x, H, W)

    x = self.norms[i](x)                     # stage-output LayerNorm

    if i == 0 and gather_indices is not None:
        x = TokenMasking.restore_tokens(x, gather_indices, N, C)  # (B, N, C_1)

    B_cur, N_cur, C_cur = x.shape
    x = x.transpose(1, 2).reshape(B_cur, C_cur, H, W)  # (B, C_i, H_i, W_i)
    features.append(x)
\end{lstlisting}
Six steps: \textbf{embed, (maybe hide), blocks, normalise, (maybe restore), reshape}. The two
``maybe'' steps are the subject of \cref{ch:masking}.

The last line does double duty. The reshaped tensor is both this stage's feature map, saved in
\code{features} for the decoder, and the next stage's input, because the loop variable \code{x} carries
it into the next iteration. Stages 2--4 therefore never see the compressed picture: stage 2 reads F1,
stage 3 reads F2, stage 4 reads F3.

The reshape at the end is the sequence-to-spatial conversion of \cref{ch:shapes}: \lstinline|transpose(1, 2)|
swaps the token and channel axes, then \code{reshape} folds the $N$ tokens back into $H \times W$ rows
and columns. The transpose is necessary because the memory order differs between the layouts.

\section{Hiding happens at stage 1 only}\label{sec:encoder:holes}

Masked training (\cref{ch:recon}) hides some tokens from the encoder so it must predict them. A natural
first guess is that every stage should hide tokens. It does not, and the source says why:
\begin{quote}\small
``Token removal only happens at Stage 1. After scattering back and reshaping to 2D, Stage 2's OPE
(Conv2d stride=2) pools over 3x3 regions, so most Stage 2 tokens contain at least some real information
even with 50\% masking. Stages 2-4 process ALL tokens -- the information loss from masking naturally
dilutes through the spatial reduction.'' (\file{seg_former_encoder.py}, lines 22--26)
\end{quote}
In plainer terms: \textbf{the holes are punched once, at the finest resolution, and then the pyramid
blurs them out by itself.} Hiding again at stage 2 would starve the coarse stages, which are exactly
where the model learns its ``what usually goes here'' knowledge.

``Most stage-2 tokens'' is a claim you can check with numbers.

\subsection{Which stage-1 tokens a stage-2 token reads}

A stage-2 token at row $h$ is computed by a 3 $\times$ 3 window, stride 2, padding 1, over the 32
$\times$ 32 grid of F1. Its window covers F1 rows $2h - 1$, $2h$, $2h + 1$. Row $-1$ is padding (a
zero that is not a token).
\begin{itemize}
\item For $h = 0$: rows $-1, 0, 1$, so only rows 0 and 1 are real: 2 rows.
\item For $h = 15$: rows $29, 30, 31$, all real: 3 rows. The padding at the bottom is never reached,
because $2 \cdot 15 + 1 = 31$ is the last row.
\end{itemize}
The same holds for columns. So in the 16 $\times$ 16 grid of stage 2:
\begin{itemize}
\item 1 token (row 0, column 0) reads $2 \times 2 = 4$ stage-1 tokens;
\item $15 + 15 = 30$ tokens (row 0 or column 0, not both) read $2 \times 3 = 6$;
\item $15 \times 15 = 225$ tokens read $3 \times 3 = 9$.
\end{itemize}
Check: $1 + 30 + 225 = 256$. Every stage-2 token is accounted for.

\subsection{The expected number of empty stage-2 tokens}

\begin{definition}[Expected value, linearity]
The \index{expected value}\emph{expected value} $\E[X]$ of a random count $X$ is its long-run average
over many repetitions. If $X = X_1 + X_2 + \dots + X_n$ is a sum of counts, then
$\E[X] = \E[X_1] + \dots + \E[X_n]$, whether or not the $X_k$ are independent
(\index{linearity of expectation}\emph{linearity of expectation}). For an indicator $X_k$ (1 if an event
happens, 0 if not), $\E[X_k]$ is the probability of the event.
\end{definition}

Call a stage-2 token \emph{empty} if every stage-1 token in its window was hidden. Treat each stage-1
token as hidden with probability $\rho = 0.65$, independently (an approximation; the exact draw hides a
fixed count, see below). A window of $m$ tokens is empty with probability $\rho^m$:
\begin{align*}
\rho^4 &= 0.65^4 = 0.65^2 \cdot 0.65^2 = 0.4225 \times 0.4225 \approx 0.1785,\\
\rho^6 &= 0.65^4 \cdot 0.65^2 \approx 0.1785 \times 0.4225 \approx 0.0754,\\
\rho^9 &= 0.65^6 \cdot 0.65^3 \approx 0.0754 \times 0.2746 \approx 0.0207 .
\end{align*}
By linearity, the expected number of empty stage-2 tokens is the sum of the 256 probabilities:
\begin{align*}
\E[\text{empty}] &= 1 \times 0.1785 + 30 \times 0.0754 + 225 \times 0.0207\\
                 &\approx 0.1785 + 2.262 + 4.66\\
                 &\approx 7.1 .
\end{align*}
About 7 of the 256 stage-2 tokens\tagC, under 3\%, have no visible data in their window. The other 97\%
see at least one survivor. The exact draw (665 of 1024 hidden, without replacement) gives $6.94$; the
script \file{research/textbook/tools/compute_masking.py} computes both\tagC. By stage 3, whose tokens
see 3 $\times$ 3 windows of stage-2 tokens, an empty token would need nine empty stage-2 tokens in a
row and column block, which essentially never happens.

\begin{keypoint}
Masking at stage 1 at 65\% leaves only about 7 of 256 stage-2 tokens without any visible input; the
pyramid fills the holes for free.
\end{keypoint}

\section{The complete masked walk-through}\label{sec:encoder:walk}

Take two fully valid patches ($B = 2$) of 128 $\times$ 128 and the training ratio $\rho = 0.65$.
The ratio comes from the training configuration
\tagM{configs/hyperspectral_segformer_exp_2.json}.
How many tokens survive? The hiding count is
$\lfloor 1024 \times 0.65 \rfloor = \lfloor 665.6 \rfloor = 665$ (\cref{ch:masking} explains the
rounding down), so $1024 - 665 = 359$ tokens are kept.

\begin{table}[htbp]\centering
\caption{One masked pass through the encoder, $B = 2$, fully valid patches, $\rho = 0.65$. Shapes were
confirmed by running the real encoder on CPU
(\texttt{research/textbook/tools/compute\_params.py}).}\label{tab:encoder:walk}
\begin{tabular}{@{}llll@{}}\toprule
Stage & Step & Shape after the step & Note \\\midrule
1 & patch embed $k4\,s4\,p0$ & \shape{(2, 1024, 32)} & all tokens \\
  & remove tokens & \shape{(2, 359, 32)} & 665 hidden tokens deleted \\
  & 2 blocks & \shape{(2, 359, 32)} & full-attention fallback \\
  & LayerNorm & \shape{(2, 359, 32)} & \\
  & restore tokens & \shape{(2, 1024, 32)} & zeros at hidden slots \\
  & reshape & \shape{(2, 32, 32, 32)} = F1 & a grid with holes \\\midrule
2 & patch embed $k3\,s2\,p1$ & \shape{(2, 256, 64)} & pools over the holes \\
  & 2 blocks, LayerNorm & \shape{(2, 256, 64)} & $R = 4$ \\
  & reshape & \shape{(2, 64, 16, 16)} = F2 & no holes \\\midrule
3 & patch embed, blocks, norm & \shape{(2, 64, 160)} & $R = 2$ \\
  & reshape & \shape{(2, 160, 8, 8)} = F3 & \\\midrule
4 & patch embed, blocks, norm & \shape{(2, 16, 256)} & $R = 1$, true full attention \\
  & reshape & \shape{(2, 256, 4, 4)} = F4 & \\\bottomrule
\end{tabular}
\end{table}

\begin{sourcenote}[: 359, not 358]
The course writes ``\textasciitilde 358'' kept tokens at $\rho = 0.65$, which is $35\%$ of $1024$ ($358.4$).
The code rounds the \emph{hidden} count down, \lstinline|(num_valid * mask_ratio).long()|
(\file{app/foundation_models/components/token_masking.py}, line 186), so a fully valid patch hides 665
and keeps 359. Running the encoder confirms 359 per item
(\file{research/textbook/tools/compute_params_output.txt}).
\end{sourcenote}

\section{Three things that look like bugs and are not}\label{sec:encoder:notbugs}

\subsection{$H$ and $W$ keep describing the full grid}

After removal there are 359 tokens, but the loop still passes $H = W = 32$ to the blocks. That is
correct and necessary. The restore step needs the full size to rebuild the grid, and the final reshape
needs it too. Meanwhile the mismatch $N \ne H \times W$ is \emph{how} ESA and Mix-FFN detect that
masking is active and take their fallback paths.

\subsection{Hidden positions become exactly zero}

The masked-autoencoder paper \citep{he2022masked} inserts a learned \index{mask token}\emph{mask token}
at hidden positions: one trained vector that means ``something was here''. This encoder writes zeros
(\code{restore_tokens} starts from \code{torch.zeros}, \file{token_masking.py}, line 402). It works
because the decoder's job here is spatial fusion of four maps rather than sequence modelling: it
fills the gaps from the surviving stage-1 tokens around them and from the fully populated F2--F4.

\subsection{F1 keeps its holes all the way to the decoder}

Nothing fills them at stage 1. F1 arrives at the decoder with zero vectors at every hidden token. That
is fine: the decoder receives four maps covering the same ground and only one of them has holes.

\section{What the encoder returns}\label{sec:encoder:returns}

\begin{lstlisting}
return features       # [F1, F2, F3, F4]
# F1: (B, 32, 32, 32)   F2: (B, 64, 16, 16)   F3: (B, 160, 8, 8)   F4: (B, 256, 4, 4)
\end{lstlisting}
A Python list of four tensors at four resolutions and four widths: the complete interface between
encoder and decoder. Any patch side that divides by 32 works, because the side is divided by 4 once and
by 2 three times ($4 \times 2 \times 2 \times 2 = 32$); for a 256 $\times$ 256 patch the grids are
$64, 32, 16, 8$.

\begin{confusions}
\confusion{Do stages 2--4 see the original image?}{No. Stage 2 reads F1, stage 3 reads F2, stage 4
reads F3. Only stage 1 touches the compressed picture.}
\confusion{Is the encoder a bottleneck like an ordinary autoencoder's?}{A hierarchical one. There is no
single narrow layer; there are four maps, and the coarsest (16 tokens of 256 numbers, 4,096 values) is
very compressed. All four go to the decoder, so the fine map is not squeezed through the coarse one.}
\confusion{Why does the width go up as the grid goes down?}{To shift capacity from ``where'' to ``what''
while each map still shrinks overall (32,768 to 4,096 numbers). It is the standard convolutional
backbone pattern, inherited by SegFormer.}
\confusion{Stage 1 has R = 8, so its attention is always reduced?}{Only without masking. In a masked
pass the token count is not $H \times W$, ESA cannot rebuild the grid, and stage 1 runs full attention
among the kept tokens (\cref{sec:encoder:attncost}).}
\confusion{Does every stage-2 token mix hidden and visible tokens?}{Most do. About 7 of 256 are
expected to see only hidden tokens at $\rho = 0.65$; they carry only bias, LayerNorm and padding
values, and the decoder must rely on F3 and F4 there.}
\confusion{Could I use a patch size other than 128?}{Yes, if it divides by 32. The token count, the
key count and the run time change; the parameter count does not (\cref{ch:params}).}
\end{confusions}

\begin{checkyourself}
\item\label{ex:encoder:1} Use \cref{eq:encoder:convout} to find the four grid sizes for a 256
$\times$ 256 patch, and the token count of each stage.
\item\label{ex:encoder:2} For a 128 $\times$ 128 patch, how many keys does each stage's attention use
in an unmasked pass, and how many score entries per head does stage 2 compute?
\item\label{ex:encoder:3} Using \cref{eq:encoder:rf}, find the receptive field and jump of a stage-2
token if stage 1 were built with kernel 7, stride 4 (the original SegFormer). Why would that break
masked training?
\item\label{ex:encoder:4} A batch is masked at $\rho = 0.5$ instead of 0.65. Using the independence
approximation, compute the expected number of empty stage-2 tokens for a 128 $\times$ 128 patch.
\item\label{ex:encoder:5} Name the six steps of one stage, and say which two only happen at stage 1 and
why the reshape at the end is used twice.
\end{checkyourself}

\nextchapter{ch:masking}
